% !PRINTED 110
\weberpage{110}
Mit Rücksicht auf die Formel Bd. I, § 145, (3) kann man dieser Relation noch die Form geben:

\weq{24}{2^{\frac{n-1}{2}} \prod_{1,\,\frac{n-1}{2}}^{\nu} \vartheta_{00}\left(\frac{2\nu}{n}\right)\vartheta_{10}\left(\frac{2\nu}{n}\right)\vartheta_{01}\left(\frac{2\nu}{n}\right) = (-1)^{\frac{n^2-1}{8}} \vartheta_{00}^{\frac{n-1}{2}}\vartheta_{10}^{\frac{n-1}{2}}\vartheta_{01}^{\frac{n-1}{2}}.}

\begin{center}
\textbf{\S\,33. \gesperrt{Die Haupttransformationen ungerader Ordnung.}}
\end{center}

Die zuletzt bewiesene Formel ist uns von Nutzen bei der Durchführung der Transformation ungerader Ordnung $n$. Wir betrachten zunächst die \gesperrt{erste Haupttransformation}. Nach § 27, (12) ist

\weq{1}{\vartheta_{11}(nu,\,n\omega)}

eine $\Theta_{11}$-Funktion $n$ter Ordnung von $u$ und $\omega$, und diese läßt sich aus ihren Nullpunkten leicht bilden, deren es im Periodenparallelogramm $n$ gibt. Die Nullpunkte von (1) sind die Werte
\[
u = \frac{\nu + \mu n\omega}{n} = \frac{\nu}{n} + \mu\omega,
\]
wenn $\nu$ und $\mu$ ganze Zahlen sind, und man erhält alle inkongruenten unter diesen Werten, wenn man $\mu$ festhält und $\nu$ ein volles Restsystem nach dem Modul $n$ durchlaufen läßt. Wir wählen das Restsystem
\[
0,\ \pm 1,\ \pm 2,\ \ldots,\ \pm\frac{(n-1)}{2},
\]
und erhalten demnach, wenn $C$ einen von $u$ unabhängigen Faktor bedeutet,

\weq{2}{C\vartheta_{11}(nu,\,n\omega) = \vartheta_{11}(u) \prod_{1,\,\frac{n-1}{2}}^{\nu} \vartheta_{11}\left(\frac{\nu}{n} + u\right)\vartheta_{11}\left(\frac{\nu}{n} - u\right),}

eine Formel, die sich nach § 22, (4) in folgender Weise auch durch die Funktionen $\vartheta_{11}(u)$, $\vartheta_{01}(u)$ ausdrücken läßt:
\[
C\vartheta_{01}^{n-1}\vartheta_{11}(nu,\,n\omega) = \vartheta_{11}(u) \prod_{1,\,\frac{n-1}{2}}^{\nu} \left[\vartheta_{11}^2\left(\frac{\nu}{n}\right)\vartheta_{01}^2(u) - \vartheta_{01}^2\left(\frac{\nu}{n}\right)\vartheta_{11}^2(u)\right].
\]
Wir ersetzen $u$ durch
\[
u + \frac{1}{2}, \quad u + \frac{\omega}{2}, \quad u + \frac{1+\omega}{2}
\]
%CONT>>
